% Options for packages loaded elsewhere
\PassOptionsToPackage{unicode}{hyperref}
\PassOptionsToPackage{hyphens}{url}
\PassOptionsToPackage{dvipsnames,svgnames,x11names}{xcolor}
\documentclass[
  american,
  11pt,
]{article}
\usepackage{xcolor}
\usepackage[margin=1in]{geometry}
\usepackage{amsmath,amssymb}
\setcounter{secnumdepth}{-\maxdimen} % remove section numbering
\usepackage{iftex}
\ifPDFTeX
  \usepackage[T1]{fontenc}
  \usepackage[utf8]{inputenc}
  \usepackage{textcomp} % provide euro and other symbols
\else % if luatex or xetex
  \usepackage{unicode-math} % this also loads fontspec
  \defaultfontfeatures{Scale=MatchLowercase}
  \defaultfontfeatures[\rmfamily]{Ligatures=TeX,Scale=1}
\fi
\usepackage{lmodern}
\ifPDFTeX\else
  % xetex/luatex font selection
\fi
% Use upquote if available, for straight quotes in verbatim environments
\IfFileExists{upquote.sty}{\usepackage{upquote}}{}
\IfFileExists{microtype.sty}{% use microtype if available
  \usepackage[]{microtype}
  \UseMicrotypeSet[protrusion]{basicmath} % disable protrusion for tt fonts
}{}
\makeatletter
\@ifundefined{KOMAClassName}{% if non-KOMA class
  \IfFileExists{parskip.sty}{%
    \usepackage{parskip}
  }{% else
    \setlength{\parindent}{0pt}
    \setlength{\parskip}{6pt plus 2pt minus 1pt}}
}{% if KOMA class
  \KOMAoptions{parskip=half}}
\makeatother
\usepackage{color}
\usepackage{fancyvrb}
\newcommand{\VerbBar}{|}
\newcommand{\VERB}{\Verb[commandchars=\\\{\}]}
\DefineVerbatimEnvironment{Highlighting}{Verbatim}{commandchars=\\\{\}}
% Add ',fontsize=\small' for more characters per line
\newenvironment{Shaded}{}{}
\newcommand{\AlertTok}[1]{\textcolor[rgb]{1.00,0.00,0.00}{\textbf{#1}}}
\newcommand{\AnnotationTok}[1]{\textcolor[rgb]{0.38,0.63,0.69}{\textbf{\textit{#1}}}}
\newcommand{\AttributeTok}[1]{\textcolor[rgb]{0.49,0.56,0.16}{#1}}
\newcommand{\BaseNTok}[1]{\textcolor[rgb]{0.25,0.63,0.44}{#1}}
\newcommand{\BuiltInTok}[1]{\textcolor[rgb]{0.00,0.50,0.00}{#1}}
\newcommand{\CharTok}[1]{\textcolor[rgb]{0.25,0.44,0.63}{#1}}
\newcommand{\CommentTok}[1]{\textcolor[rgb]{0.38,0.63,0.69}{\textit{#1}}}
\newcommand{\CommentVarTok}[1]{\textcolor[rgb]{0.38,0.63,0.69}{\textbf{\textit{#1}}}}
\newcommand{\ConstantTok}[1]{\textcolor[rgb]{0.53,0.00,0.00}{#1}}
\newcommand{\ControlFlowTok}[1]{\textcolor[rgb]{0.00,0.44,0.13}{\textbf{#1}}}
\newcommand{\DataTypeTok}[1]{\textcolor[rgb]{0.56,0.13,0.00}{#1}}
\newcommand{\DecValTok}[1]{\textcolor[rgb]{0.25,0.63,0.44}{#1}}
\newcommand{\DocumentationTok}[1]{\textcolor[rgb]{0.73,0.13,0.13}{\textit{#1}}}
\newcommand{\ErrorTok}[1]{\textcolor[rgb]{1.00,0.00,0.00}{\textbf{#1}}}
\newcommand{\ExtensionTok}[1]{#1}
\newcommand{\FloatTok}[1]{\textcolor[rgb]{0.25,0.63,0.44}{#1}}
\newcommand{\FunctionTok}[1]{\textcolor[rgb]{0.02,0.16,0.49}{#1}}
\newcommand{\ImportTok}[1]{\textcolor[rgb]{0.00,0.50,0.00}{\textbf{#1}}}
\newcommand{\InformationTok}[1]{\textcolor[rgb]{0.38,0.63,0.69}{\textbf{\textit{#1}}}}
\newcommand{\KeywordTok}[1]{\textcolor[rgb]{0.00,0.44,0.13}{\textbf{#1}}}
\newcommand{\NormalTok}[1]{#1}
\newcommand{\OperatorTok}[1]{\textcolor[rgb]{0.40,0.40,0.40}{#1}}
\newcommand{\OtherTok}[1]{\textcolor[rgb]{0.00,0.44,0.13}{#1}}
\newcommand{\PreprocessorTok}[1]{\textcolor[rgb]{0.74,0.48,0.00}{#1}}
\newcommand{\RegionMarkerTok}[1]{#1}
\newcommand{\SpecialCharTok}[1]{\textcolor[rgb]{0.25,0.44,0.63}{#1}}
\newcommand{\SpecialStringTok}[1]{\textcolor[rgb]{0.73,0.40,0.53}{#1}}
\newcommand{\StringTok}[1]{\textcolor[rgb]{0.25,0.44,0.63}{#1}}
\newcommand{\VariableTok}[1]{\textcolor[rgb]{0.10,0.09,0.49}{#1}}
\newcommand{\VerbatimStringTok}[1]{\textcolor[rgb]{0.25,0.44,0.63}{#1}}
\newcommand{\WarningTok}[1]{\textcolor[rgb]{0.38,0.63,0.69}{\textbf{\textit{#1}}}}
\usepackage{longtable,booktabs,array}
\newcounter{none} % for unnumbered tables
\usepackage{calc} % for calculating minipage widths
% Correct order of tables after \paragraph or \subparagraph
\usepackage{etoolbox}
\makeatletter
\patchcmd\longtable{\par}{\if@noskipsec\mbox{}\fi\par}{}{}
\makeatother
% Allow footnotes in longtable head/foot
\IfFileExists{footnotehyper.sty}{\usepackage{footnotehyper}}{\usepackage{footnote}}
\makesavenoteenv{longtable}
\ifLuaTeX
\usepackage[bidi=basic,shorthands=off]{babel}
\else
\usepackage[bidi=default,shorthands=off]{babel}
\fi
\ifLuaTeX
  \usepackage{selnolig} % disable illegal ligatures
\fi
\setlength{\emergencystretch}{3em} % prevent overfull lines
\providecommand{\tightlist}{%
  \setlength{\itemsep}{0pt}\setlength{\parskip}{0pt}}
\usepackage[dvipsnames]{xcolor}
\usepackage{microtype}
\usepackage{fancyhdr}
\pagestyle{fancy}
\fancyhf{}
\fancyhead[L]{\small Portfolio Risk Engine}
\fancyhead[R]{\small Technical Paper}
\fancyfoot[C]{\thepage}
\setlength{\headheight}{14pt}
\setlength{\emergencystretch}{3em}
\usepackage{bookmark}
\IfFileExists{xurl.sty}{\usepackage{xurl}}{} % add URL line breaks if available
\urlstyle{same}
\hypersetup{
  pdftitle={Portfolio Risk{,} Sensitivity{,} and Constrained Allocation},
  pdfauthor={Eshan Garkhedkar},
  pdflang={en-US},
  colorlinks=true,
  linkcolor={NavyBlue},
  filecolor={Maroon},
  citecolor={Blue},
  urlcolor={NavyBlue},
  pdfcreator={LaTeX via pandoc}}

\title{Portfolio Risk, Sensitivity, and Constrained Allocation}
\usepackage{etoolbox}
\makeatletter
\providecommand{\subtitle}[1]{% add subtitle to \maketitle
  \apptocmd{\@title}{\par {\large #1 \par}}{}{}
}
\makeatother
\subtitle{Mathematical foundations, C++ implementation, and independent
validation}
\author{Eshan Garkhedkar}
\date{}

\begin{document}
\maketitle
\begin{abstract}
This project develops an auditable numerical engine for portfolio risk,
attribution, sensitivity, and contribution-only allocation. We derive
the mathematical models, implement them in C++ with a scoped
reverse-mode differentiation tape, and retain independent Python
reference calculations. The central distinction is between arithmetic
correctness, optimization correctness, and the empirical adequacy of
estimated covariance. The paper states the domains of the derivatives,
formulates the constrained quadratic problem, documents data
assumptions, and separates measured validation from unperformed
empirical work. Historical covariance is not treated as a known law of
future returns. This is a research and educational software project.
\end{abstract}

{
\setcounter{tocdepth}{2}
\tableofcontents
}
\section{1. Definitions, prerequisites, and risk
estimation}\label{definitions-prerequisites-and-risk-estimation}

\subsection{Practical question}\label{practical-question}

\textbf{Symbols used in this subsection}
\begin{itemize}
\item \(i,t\): Asset index and observation time; $t-1$ is the preceding observation.
\item \(r_{t,i},r_t\): Simple return for asset $i$ at time $t$, and vector of returns at that time.
\item \(P^*_{t,i}\): Consistently adjusted historical price; the star marks an adjustment label.
\item \(n,\mathbb R^n\): Number of assets and the space of real vectors with $n$ entries.
\item \(h_i,V,w_i,w\): Dollar holding, total value, individual portfolio weight, and vector of weights.
\item \(\mathbf1,T\): All-ones vector, and transpose when used as a superscript.
\item \(X,Y\): Random variables in the expectation, variance, and covariance definitions.
\item \(E,\operatorname{Var},\operatorname{Cov}\): Expectation, variance, and covariance operators.
\item \(\sum,\ge,>\): Summation, greater-than-or-equal comparison, and strict greater-than comparison.
\end{itemize}


How much dispersion does a fixed portfolio have under an estimated
covariance model, and which assumptions enter that estimate?

For asset \(i\), a simple one-period return is
\(r_{t,i}=P^*_{t,i}/P^*_{t-1,i}-1\). Here \(P^*\) denotes a positive,
consistently adjusted historical price series whose economic
interpretation has been checked. This notation alone does not establish
that dividends are reinvested. Current holdings require a separate
unadjusted valuation price.

Let \(r_t\in\mathbb R^n\), \(w\in\mathbb R^n\), \(w_i\ge0\) and
\(\mathbf1^Tw=1\). Returns and weights are dimensionless. If \(h_i\) are
dollar holdings, \(V=\sum_i h_i>0\) and \(w_i=h_i/V\). Covariance is in
squared-return units for a declared period; volatility is in return
units for the same period.

Expectation is a probability-weighted average. Variance
\(\operatorname{Var}(X)=E[(X-E[X])^2]\) measures squared deviations
around it. Covariance \(\operatorname{Cov}(X,Y)=E[(X-E[X])(Y-E[Y])]\)
measures joint deviations. Correlation divides covariance by the product
of nonzero standard deviations. A zero-variance variable has no defined
correlation coefficient under this formula.

\subsection{Portfolio quadratic form}\label{portfolio-quadratic-form}

\textbf{Symbols used in this subsection}
\begin{itemize}
\item \(r,r_p\): Asset-return vector and resulting portfolio return.
\item \(w,w_i,w_j\): Weight vector and its entries at asset indices $i,j$.
\item \(E,\operatorname{Var},\sum_i,\sum_j\): Expectation, variance, and sums across asset indices.
\item \(a\): Arbitrary real vector used to test nonnegative quadratic variance.
\item \(T\): Transpose of a vector or matrix; not the observation count in this subsection.
\item \(Q,\Lambda\): Orthogonal eigenvector matrix and diagonal eigenvalue matrix.
\item \(\lambda_{max},\lambda_{min}\): Largest and smallest eigenvalues; subscripts indicate extrema.
\end{itemize}


With fixed weights, \(r_p=w^Tr\). Expanding the square gives

\[\operatorname{Var}(r_p)=E\left[\left(\sum_i w_i(r_i-E[r_i])\right)^2\right]
=\sum_i\sum_jw_iw_j\Sigma_{ij}=w^T\Sigma w.\]

Thus \(\sigma(w)=\sqrt{w^T\Sigma w}\). Covariance is symmetric and
positive semidefinite (PSD):
\(a^T\Sigma a=\operatorname{Var}(a^Tr)\ge0\) for every real \(a\). The
spectral theorem writes \(\Sigma=Q\Lambda Q^T\); PSD means nonnegative
eigenvalues. Eigenvalues measure variance along orthogonal directions,
not individual asset risks.

Singularity means some direction has zero variance under this model. For
positive-definite \(\Sigma\), the spectral condition number is
\(\lambda_{max}/\lambda_{min}\); a large ratio warns that inversion can
amplify perturbations. Risk evaluation needs matrix multiplication, not
inversion. Singularity is therefore valid and must not be rejected
simply because a Cholesky factorization fails. See
\href{https://docs.mosek.com/portfolio-cookbook/markowitz.html}{MOSEK's
portfolio formulation}.

\subsection{Sample estimator and
normalization}\label{sample-estimator-and-normalization}

\textbf{Symbols used in this subsection}
\begin{itemize}
\item \(T,n,t\): Number of return observations, number of assets, and observation index.
\item \(r_t,\bar r\): Return vector at observation $t$ and its sample mean.
\item \(X\): Centered return matrix, with observations as rows and assets as columns.
\item \(S\): Sample covariance matrix using the $T-1$ denominator.
\item \(a\): Arbitrary test vector for the quadratic form.
\item \({}^T,\|\cdot\|,\operatorname{rank},\min\): Transpose, Euclidean norm, matrix rank, and minimum operator.
\item \(\mathbb R^{T\times n},\sum_t\): Real matrices with $T$ rows and $n$ columns, and sum over observations.
\end{itemize}


For \(T\ge2\) aligned return rows, \(\bar r=T^{-1}\sum_t r_t\) and
centered matrix \(X\in\mathbb R^{T\times n}\):

\[S=\frac{1}{T-1}X^TX.\]

This uses the sample normalization; under independent identically
distributed rows with finite covariance it is unbiased. Rows of market
data need not meet those assumptions. Since \(a^TSa=\|Xa\|^2/(T-1)\),
this construction preserves PSD. Centering implies
\(\operatorname{rank}(S)\le\min(n,T-1)\). Pairwise deletion with
different date sets need not preserve PSD; the product must use a common
set of complete daily returns.

Example with three return rows \((.01,.02),(-.01,0),(0,-.02)\) has zero
mean and

\[S=\begin{pmatrix}.0001&.0001\\.0001&.0004\end{pmatrix}.\]

The reference independently matches this result and NumPy's
\texttt{cov(...,\ ddof=1)}.

\subsection{Shrinkage and its precise
status}\label{shrinkage-and-its-precise-status}

\textbf{Symbols used in this subsection}
\begin{itemize}
\item \(S,S_{ML},S_\alpha\): Sample covariance, empirical covariance using denominator $T$, and shrunk covariance.
\item \(\alpha,\hat\alpha\): Chosen shrinkage coefficient and fitted shrinkage coefficient; a hat marks estimation.
\item \(\mu,n,I\): Mean diagonal variance, asset count, and identity matrix.
\item \(\operatorname{tr},\lambda_i\): Sum of diagonal entries and eigenvalue indexed by $i$.
\item \(X,T,{}^T\): Centered return matrix, observation count, and transpose.
\end{itemize}


For a PSD estimate \(S\), define \(\mu=\operatorname{tr}(S)/n\) and

\[S_\alpha=(1-\alpha)S+\alpha\mu I,\qquad 0\le\alpha\le1.\]

It is PSD as a convex combination of PSD matrices. If \(\alpha>0\) and
\(\mu>0\), all eigenvalues \((1-\alpha)\lambda_i+\alpha\mu\) are
positive. This improves conditioning but introduces a structured bias
toward equal variances and zero covariance; it is not a guarantee of
better future allocation.

The engine implements a fixed-\(\alpha\) sensitivity path. Separately,
Python fits the Ledoit--Wolf scaled-identity estimator through
scikit-learn and records the estimated coefficient. Its covariance-loss
criterion is not realized portfolio returns. See the
\href{https://www.sciencedirect.com/science/article/pii/S0047259X03000964}{original
paper} and
\href{https://scikit-learn.org/stable/modules/generated/sklearn.covariance.LedoitWolf.html}{estimator
documentation}.

The fitted estimator uses the centered empirical covariance
\(S_{ML}=X^TX/T\), while the sample option uses \(S=X^TX/(T-1)\). Tests
reconstruct the fitted result as
\((1-\hat\alpha)S_{ML}+\hat\alpha\operatorname{tr}(S_{ML})I/n\). Both
outputs are multiplied by the declared annualization factor. Thus
estimator comparisons include a normalization difference as well as
shrinkage. A fixed arbitrary coefficient is never labeled a fitted
Ledoit--Wolf estimate.

\subsection{Annualization and
interpretation}\label{annualization-and-interpretation}

\textbf{Symbols used in this subsection}
\begin{itemize}
\item \(K,d\): Number of additive periods and a subscript identifying daily quantities.
\end{itemize}


For an additive \(K\)-period return sum, covariance equals \(K\Sigma_d\)
only if cross-period covariance terms vanish and period covariances are
identical. With serial dependence the omitted cross terms matter. Annual
volatility \(\sqrt{252}\sigma_d\) is a convention under these
assumptions, not the exact variance of compounded annual simple returns.
Display both period and multiplier. The example in chapter 2 is already
annual and must not be annualized again.

This is a constant-weight risk estimate. A buy-and-hold historical
account has drifting weights; reproducing that path requires holdings
evolution. Historical covariance estimates future risk only through a
modeling assumption of relevant history. The interface compares windows
of 63, 126 and 252 completed common daily returns when available. It
labels their spread sensitivity rather than calibrated uncertainty.

Implementation:
\href{https://github.com/eshangark/portfolio-risk-engine/blob/main/src/portfolio_risk/reference.py}{reference.py}
and
\href{https://github.com/eshangark/portfolio-risk-engine/blob/main/src/portfolio_risk/estimation.py}{estimation.py}.
Validation includes hand values, scaling, permutation, PSD rejection,
singular cases and estimator normalization. Exchange-calendar processing
is implemented and tested with synthetic dates; empirical estimation
quality is explored in chapter 5, without a claim of general predictive validity.

\section{2. Differentiation and a worked
example}\label{differentiation-and-a-worked-example}

\subsection{From calculus to
decisions}\label{from-calculus-to-decisions}

\textbf{Symbols used in this subsection}
\begin{itemize}
\item \(w,v(w)\): Weight vector and variance evaluated at those weights.
\item \(\delta\): Small vector perturbation to the weights.
\item \(g,\nabla,\nabla_w\): Volatility gradient, gradient operator, and differentiation with respect to weights.
\item \(d,\mathbf1\): Feasible transfer direction and all-ones vector.
\item \(e_A,e_B\): Unit vectors selecting assets A and B.
\item \(\epsilon,g_A,g_B\): Small transfer size and gradient entries for assets A and B.
\item \({}^T,O(\epsilon^2)\): Transpose and an error whose magnitude is bounded by a constant times the squared transfer size near zero.
\end{itemize}


For a fixed symmetric \(\Sigma\), expand
\((w+\delta)^T\Sigma(w+\delta)\) to obtain

\[v(w+\delta)=v(w)+2\delta^T\Sigma w+\delta^T\Sigma\delta.\]

Therefore \(\nabla v=2\Sigma w\). Applying the scalar chain rule to
\(\sqrt v\) yields

\[g=\nabla_w\sigma=\frac{\Sigma w}{\sigma},\qquad \sigma>0.\]

The derivative is with respect to unconstrained coordinates even though
valid portfolios lie on a simplex. A feasible first-order direction
\(d\) must satisfy \(\mathbf1^Td=0\) and remain inside nonnegativity
bounds. A transfer from A to B has \(d=e_B-e_A\), so

\[\sigma(w+\epsilon d)-\sigma(w)=\epsilon(g_B-g_A)+O(\epsilon^2).\]

At a boundary a feasible direction may allow only one sign of
\(\epsilon\). Do not confuse this derivative with increasing B while
holding all other weights fixed.

\subsection{Risk attribution}\label{risk-attribution}

\textbf{Symbols used in this subsection}
\begin{itemize}
\item \(w,w_i,g\): Weight vector, asset-$i$ weight, and volatility-gradient vector.
\item \(RC_i,i,\sum_i\): Asset-$i$ volatility contribution, asset index, and sum across assets.
\item \({}^T,(\cdot)_i\): Transpose and extraction of vector entry $i$.
\end{itemize}


Volatility is homogeneous of degree one for positive scalings. Euler's
identity gives \(w^Tg=\sigma\); directly,

\[\sum_i RC_i=\sum_i w_i\frac{(\Sigma w)_i}{\sigma}
=\frac{w^T\Sigma w}{\sigma}=\sigma.\]

\(RC_i\) has the same units as volatility. Capital share \(w_i\) and
risk share \(RC_i/\sigma\) answer different questions. Contributions can
be negative, even in a long-only portfolio, if covariance gives a
sufficiently strong hedge; they are an attribution, not independent
stand-alone risks. The attribution framework is described in
\href{https://docs.mosek.com/portfolio-cookbook/risk_parity.html}{MOSEK's
risk budgeting chapter}.

\subsection{New cash changes the
denominator}\label{new-cash-changes-the-denominator}

\textbf{Symbols used in this subsection}
\begin{itemize}
\item \(h,h_i,H,V\): Dollar holdings vector, holding of asset $i$, and two labels for existing total wealth.
\item \(w,w_i,g,g_j\): Weight vector, weight $i$, volatility-gradient vector, and gradient entry $j$.
\item \(i,j,\delta_{ij}\): Asset indices and Kronecker delta: one for equal indices, zero otherwise.
\item \(\partial,\nabla_h,{}^T,\sum_i\): Partial derivative, gradient with respect to holdings, transpose, and sum over assets.
\item \(b,B,x\): Transfer amount, new contribution budget, and purchase-dollar vector.
\item \(g_A,g_B\): Gradient entries for the two assets used in the transfer example.
\end{itemize}


For dollar holdings \(h\), \(H=\sum_i h_i>0\), and \(w=h/H\):

\[\frac{\partial w_i}{\partial h_j}=\frac{\delta_{ij}-w_i}{H},\qquad
\frac{\partial\sigma}{\partial h_j}=\frac{g_j-w^Tg}{H}=\frac{g_j-\sigma}{H}.\]

The units are volatility per dollar. This derivative includes the effect
on all weights. Adding dollars proportionally to every existing holding
leaves weights unchanged: \(h^T\nabla_h\sigma=0\). A transfer of \(b\)
dollars keeps total wealth constant and instead has derivative
\(b(g_B-g_A)/H\).

For a fixed contribution \(B\) fully distributed through purchases
\(x\), the feasible total is constant: \(w(x)=(h+x)/(V+B)\). Its
objective gradient in purchase coordinates is \(2\Sigma(h+x)/(V+B)^2\).
Off the budget constraint, that fixed-denominator extension is a
computational definition, not the normalization for arbitrary additional
wealth. Both give the same directional derivative along feasible
transfers of purchases.

\subsection{Hand-worked example}\label{hand-worked-example}

\textbf{Symbols used in this subsection}
\begin{itemize}
\item \(w,v,g,RC,h\): Weights, portfolio variance, volatility gradient, risk-contribution vector, and dollar holdings.
\item \(A,B\): Labels for the first and second assets in this example; B here is not a contribution budget.
\end{itemize}


Take \(w=(.6,.4)\) and annual
\(\Sigma=\begin{pmatrix}.04&.006\\.006&.01\end{pmatrix}\). The asset
volatilities are .2 and .1; the implied correlation is
\(.006/(.2\cdot.1)=.3\).

\[\Sigma w=(.0264,.0076),\quad
v=.6^2(.04)+2(.6)(.4)(.006)+.4^2(.01)=.01888.\]

Thus \(\sigma=0.1374045123\), \(g=(.1921334282,.0553111384)\), and
\(RC=(.1152800569,.0221244554)\). A accounts for 83.8983\% of total
volatility attribution despite holding 60\% of capital.

A transfer of one percentage point of portfolio weight from A to B has
linear change \(-.0013682229\) and exact change \(-.0013611369\) in
volatility units. Multiply by 100 to express volatility percentage
points: approximately \(-.13682\) versus \(-.13611\).

Start with \(h=(600,400)\) dollars and buy \(100\) of B using new cash.
Weights become \((600/1100,500/1100)\). Exact volatility change is
\(-.0072425026\) versus the local dollar approximation \(-.0082093374\).
A 10\% wealth increase is large enough that the approximation error is
visible. The conclusion concerns percentage-return volatility; dollar
risk can increase when invested wealth grows.

\subsection{Sensitivity to a valid covariance
path}\label{sensitivity-to-a-valid-covariance-path}

\textbf{Symbols used in this subsection}
\begin{itemize}
\item \(S,F,S_\alpha\): Base covariance, target covariance, and covariance along the shrinkage path.
\item \(\mu,I,\alpha\): Mean diagonal variance, identity matrix, and shrinkage coefficient.
\item \(w,{}^T,\partial/\partial\alpha\): Weight vector, transpose, and partial differentiation with respect to shrinkage.
\end{itemize}


Keep base \(S\), target \(F=\mu I\), and weights fixed. On
\(S_\alpha=(1-\alpha)S+\alpha F\):

\[\frac{\partial\sigma}{\partial\alpha}=\frac{w^T(F-S)w}{2\sigma}.\]

Only \(\alpha\in[0,1]\) is admitted. At endpoints use admissible
one-sided perturbations for checks. This derivative does not
differentiate the estimated Ledoit--Wolf coefficient with respect to
data. Changing a single correlation arbitrarily can destroy PSD; use
model-valid paths instead.

\subsection{Implemented custom reverse
mode}\label{implemented-custom-reverse-mode}

\textbf{Symbols used in this subsection}
\begin{itemize}
\item \(u,v,z\): Two scalar inputs and their product in the tape example; $v$ here is not portfolio variance.
\item \(\bar u,\bar v,\bar z\): Adjoints: derivatives of the final output with respect to these intermediate values.
\item \(\mathrel{+}=\): Accumulate into an existing adjoint rather than overwrite it.
\item \(x,q,f,df/dx\): Scalar input, squared intermediate, final output, and derivative of the output with respect to the input.
\item \(m,n,O(\cdot)\): Operation count, asset count, and asymptotic upper-bound notation for computation or storage.
\end{itemize}


The C++ tape stores immutable node values and parent edges with local
derivatives. Each reverse pass allocates fresh zero adjoints,
initializes the output adjoint to one, and traverses backward to add
chain-rule contributions. For \(z=u v\), accumulate
\(\bar u\mathrel{+}=\bar z v\) and \(\bar v\mathrel{+}=\bar z u\).
Addition, subtraction, multiplication, division and square root are
implemented. If \(u=v\), both edge contributions accumulate. For
example, \(q=x^2\) and \(f=q+q\) produce \(df/dx=4x\), including the
shared intermediate correctly.

A tape owns node storage. Handles retain shared ownership and carry a
generation number to reject cross-tape mixing and use after reset.
Graphs are rebuilt per calculation, including every optimizer gradient
evaluation; stored local derivatives are never reused with changed
inputs. Division by zero, square root at nonpositive values and
nonfinite arithmetic fail explicitly. The risk wrapper returns zero
volatility with its usual gradient unavailable, even when a special
constant-zero model admits a derivative. With \(m\) operations and
bounded parents per node, forward evaluation and one reverse pass
require \(O(m)\) work and storage. The dense quadratic graph has
\(O(n^2)\) operations; no speed advantage over the closed form is
asserted.

AD evaluates derivatives of the implemented arithmetic graph up to
floating-point effects; it does not cure statistical model error. The
closed-form volatility gradient is simple and remains an independent
oracle. AD supplies implementation depth and handles composed
normalization and parameter paths. See
\href{https://jmlr.org/papers/volume18/17-468/17-468.pdf}{Baydin et al.,
section 3}.

Independent formulas remain in
\href{https://github.com/eshangark/portfolio-risk-engine/blob/main/src/portfolio_risk/reference.py}{reference.py}.
The C++ engine evaluates risk, normalized-dollar sensitivities and
fixed-shrinkage sensitivities through the tape. Python tests compare
these against analytical formulas and finite differences; native tests
also exercise handle lifetime, reset and shared nodes under address and
undefined-behavior sanitizers.

\section{3. Constrained allocation}\label{constrained-allocation}

\subsection{Practical question and
domain}\label{practical-question-and-domain}

\textbf{Symbols used in this subsection}
\begin{itemize}
\item \(h,V,B,D\): Existing dollar holdings, existing wealth, contribution budget, and final total wealth.
\item \(x,x_i,c_i,u_i\): Purchase vector, purchase for asset $i$, final-weight cap, and dollar purchase capacity.
\item \(n,i,\mathbf1,\mathbb R^n\): Asset count, asset index, all-ones vector, and real-vector space.
\item \(f,\min_x\): Portfolio variance objective and minimization over purchases.
\item \({}^T,\nabla,\nabla^2,\succeq0\): Transpose, gradient, Hessian (second-derivative matrix), and positive-semidefinite ordering.
\end{itemize}


Given nonnegative dollar holdings \(h\in\mathbb R^n\), wealth
\(V=\mathbf1^Th\), contribution \(B\ge0\), and \(D=V+B>0\), distribute
purchases \(x\) without selling. Assume a fixed PSD covariance,
continuous dollar amounts, full investment of \(B\) and optional
post-contribution caps \(c_i\in[0,1]\).

\[\min_x f(x)=\frac{(h+x)^T\Sigma(h+x)}{D^2},\qquad
\mathbf1^Tx=B,\quad 0\le x_i\le u_i=c_iD-h_i.\]

The objective has squared-return units; \(x\) and \(u\) are dollars. The
gradient has squared-return units per dollar:

\[\nabla f(x)=\frac{2\Sigma(h+x)}{D^2},\qquad
\nabla^2 f=\frac{2\Sigma}{D^2}\succeq0.\]

PSD curvature proves convexity: along every line the second derivative
is nonnegative. The box intersected with the budget hyperplane is convex
and compact when feasible. A continuous objective has a minimizer there.
Strict positivity of curvature on nonzero feasible directions ensures
uniqueness; singular covariance can instead produce many optimal
allocations. Weight agreement is consequently an inadequate solver
check.

\subsection{Feasibility before
optimization}\label{feasibility-before-optimization}

\textbf{Symbols used in this subsection}
\begin{itemize}
\item \(h,h_i,V,B,c,c_i,u,u_i\): Holdings vector and entry, existing wealth, budget, cap vector and entry, and purchase-capacity vector and entry.
\item \(i,\sum_i,\le,\ge\): Asset index, sum across assets, less-than-or-equal, and greater-than-or-equal comparisons.
\end{itemize}


With upper bounds only, feasibility is equivalent to \(u_i\ge0\) for all
\(i\) and \(\sum_i u_i\ge B\). Necessity follows by summing upper
bounds; sufficiency follows by filling nonnegative capacities until
\(B\) is used. An existing overweight asset may become permissible after
dilution by the contribution. Evaluate \(h_i\le c_i(V+B)\), not
\(h_i/V\le c_i\).

For \(h=(600,400)\), \(B=200\), and \(c=(.6,.6)\), \(u=(120,320)\) and
the budget is feasible. Tightening the first cap to .4 makes its
capacity \(-120\), so no purchase-only solution exists. The reference
implements these checks in
\href{https://github.com/eshangark/portfolio-risk-engine/blob/main/src/portfolio_risk/reference.py}{\texttt{purchase\_bounds}};
the implemented solver's residual tolerances are specified below.

\subsection{Algorithm and projection}\label{algorithm-and-projection}

\textbf{Symbols used in this subsection}
\begin{itemize}
\item \(x,x^k,x^{k+1},x_i\): Purchase vector, current iterate, next iterate, and purchase for asset $i$; $k$ counts updates.
\item \(C,\Pi_C\): Feasible set and Euclidean projection onto that set.
\item \(u,B,h,c,D,\mathbf1\): Purchase capacities, budget, existing holdings, caps, final wealth, and all-ones vector.
\item \(\eta,L,\lambda_{max}\): Step size, gradient Lipschitz bound, and largest eigenvalue.
\item \(f,\nabla f,z,\tau\): Variance objective, its gradient, trial point, and scalar projection threshold.
\item \(\min,\max,\sum_i,\|\cdot\|,{}^T\): Minimum, maximum, sum across assets, Euclidean norm, and transpose.
\item \(\ell,y,b,a,C_y\): Existing holdings divided by final wealth, purchase fractions, budget fraction, capacity fractions, and their feasible set.
\item \(s,Q,F\): Largest covariance eigenvalue, covariance divided by that eigenvalue, and scaled variance objective.
\item \(p,z_p,n\): Pivot asset index, pivot score, and asset count in the extreme-range projection argument.
\item \(d,y^0,y^*,\nabla F\): Iterate displacement, initial purchase fractions, optimal purchase fractions, and scaled-objective gradient; star marks an optimum.
\end{itemize}


Projected gradient is scoped, transparent, and independently testable.
Start from a feasible point and iterate

\[x^{k+1}=\Pi_C(x^k-\eta\nabla f(x^k)),\quad
C=\{x:0\le x\le u,\;\mathbf1^Tx=B\}.\]

For \(L=2\lambda_{max}(\Sigma)/D^2>0\), use \(0<\eta\le1/L\) as an
initial fixed-step design. If \(L=0\), every feasible point is optimal
and no reciprocal step is formed. Projected gradient is the
proximal-gradient method with an indicator of the feasible set; see
\href{https://web.stanford.edu/~boyd/papers/pdf/prox_algs.pdf}{Parikh
and Boyd, section 4.2}.

To project a trial point \(z\), minimize \(\frac12\|x-z\|^2\) over
\(C\). Stationarity with an equality multiplier \(\tau\) and box
clipping yields

\[x_i(\tau)=\min\{u_i,\max\{0,z_i-\tau\}\},\quad
\sum_i x_i(\tau)=B.\]

The sum is continuous and nonincreasing. For \(0<B<\sum_i u_i\), bracket
\(\tau\) between \(\min_i(z_i-u_i)\) and \(\max_i z_i\), then bisect
until the budget residual meets a documented scale-aware tolerance. The
projected point is unique, though the multiplier need not be. Handle
\(B=0\) and \(B=\sum_i u_i\) directly. Validate projection feasibility,
idempotence, and the variational inequality
\((z-\Pi_Cz)^T(y-\Pi_Cz)\le0\) for feasible \(y\).

The implemented coordinates are \(\ell=h/D\), \(y=x/D\), \(b=B/D\) and
\(a=c-\ell\). Thus \(w=\ell+y\), with feasible set
\(C_y=\{y:0\le y\le a,\ \mathbf1^Ty=b\}\).
This avoids dollar-scale conditioning and remains defined when \(B=0\).
For \(s=\lambda_{max}(\Sigma)>0\), define \(Q=\Sigma/s\) and
\(F(y)=(\ell+y)^TQ(\ell+y)\). The gradient \(2Q(\ell+y)\) has Lipschitz
constant 2, so the code takes step \(\eta=1/2\). Each projection uses
rescaled bisection for 160 iterations and a final budget-residual
correction inside the box. The baseline projects equal purchase
fractions onto the feasible set. Zero contribution or zero covariance
has an immediate feasible optimum.

An extreme-range audit exposed a weakness in scaling by the full score
range: a difference of order $10^{300}$ can erase a unit budget. The
implementation now first replaces each upper bound by $\min(u_i,B)$,
which leaves the feasible set unchanged. Sort scores in descending order
and choose pivot $p$ where their cumulative capped capacity first reaches
$B$. At threshold $z_p$, purchased capacity is below $B$; at threshold
$z_p-nB$, every asset at or above the pivot is saturated, giving capacity
at least $B$. Hence a valid threshold lies in $[z_p-nB,z_p]$.
After shifting by $z_p$ and scaling by $B$, clip scores to $[-n,n]$:
clipped extremes must be zero or saturated throughout this threshold
interval. Bisection therefore retains the budget geometry independently
of the original score range. This argument concerns exact arithmetic;
finite-precision residual correction and explicit diagnostics remain necessary.

To see why this step descends, let \(d=y^{k+1}-y^k\). Projection
optimality gives \(\nabla F(y^k)^Td\le-\|d\|^2/\eta\). The
Lipschitz-gradient bound then gives
\(F(y^{k+1})\le F(y^k)-(1/\eta-L/2)\|d\|^2\). For \(L=2\) and
\(\eta=1/2\), the coefficient is one. In exact arithmetic, convex
projected gradient also satisfies
\(F(y^k)-F(y^*)\le\|y^0-y^*\|^2/(2\eta k)\) for \(k\ge1\). This is an
objective bound, not a guarantee of unique weights. Floating-point
projection error and a finite iteration limit make explicit diagnostics
necessary.

\subsection{Optimality and independent
validation}\label{optimality-and-independent-validation}

\textbf{Symbols used in this subsection}
\begin{itemize}
\item \(f,x,u,\mathbf1\): Dollar-coordinate variance objective, purchase vector, capacities, and all-ones vector.
\item \(\nu,\lambda,\gamma\): Budget-equality multiplier, nonnegative lower-bound multipliers, and nonnegative upper-bound multipliers.
\item \(\lambda_i,\gamma_i,i\): Individual constraint multipliers and asset index; $\lambda$ here denotes multipliers, not covariance eigenvalues.
\item \(\eta,G_\eta,\Pi_C,\nabla f\): Step size, projected-gradient mapping, feasible-set projection, and objective gradient.
\item \(y,C_y,F,s,\|\cdot\|_2\): Purchase fractions, their feasible set, scaled variance objective, covariance scale, and Euclidean norm.
\end{itemize}


For the original dollar problem, KKT stationarity is

\[\nabla f+\nu\mathbf1-\lambda+\gamma=0,\]

with unrestricted \(\nu\), nonnegative \(\lambda,\gamma\), feasibility, \(\lambda_ix_i=0\) and
\(\gamma_i(x_i-u_i)=0\). For this convex quadratic on a polyhedron, an
appropriate first-order variational inequality characterizes global
optima. A convenient numerical diagnostic is

\[G_\eta(x)=\frac{x-\Pi_C(x-\eta\nabla f(x))}{\eta}.\]

It vanishes exactly at a constrained optimum. The implementation stops
when the scaled mapping norm \(2\|y-\Pi_{C_y}(y-\nabla F(y)/2)\|_2\) is at
most \(10^{-9}\) by default, or after 10,000 updates. This absolute,
dimensionless threshold is adjustable. Multiplying this norm by \(s\)
gives the mapping norm in weight coordinates for the unscaled variance
objective with step \(1/(2s)\); it is not a dollar-coordinate norm. The
output also reports dollar budget residual, maximum weight-bound
violation, iteration count, stopping reason, objective and improvement
over the feasible baseline. Constraint feasibility admits only
\(10^{-12}\) weight-scale roundoff before projection. Independent tests
compare full-rank and rank-deficient objectives with OSQP, emphasizing
objective agreement instead of unique weight agreement.

The C++ solver obtains the composed objective gradient through the tape
by default; an analytical-gradient mode is also checked. Covariance
stays fixed within each solve. The interface repeats solves across
predeclared windows and covariance estimators when enough observations
exist. Synthetic convergence and solver-agreement evidence is reported
in chapter 5, alongside a selected historical-market allocation stability study.
That study is descriptive and does not establish future stability.

Minimum estimated variance ignores expected return, investor objectives,
tax, fees, execution and estimation error. It answers this explicit
constrained mathematical question; it does not define the best personal
investment portfolio.

\section{4. Architecture and data
methodology}\label{architecture-and-data-methodology}

\subsection{Implemented separation of
responsibilities}\label{implemented-separation-of-responsibilities}

The dependency chain is provider adapter, validated returns, Python
covariance estimation, C++ risk/AD/optimization, then Streamlit
presentation. The C++17 engine accepts arrays and explicit parameters
without network or dashboard dependencies. CMake and scikit-build-core
package a \href{https://pybind11.readthedocs.io/en/stable/}{pybind11}
extension. Eigen supplies dense arithmetic and symmetric
eigendecomposition. The independent NumPy reference remains separate.

Bindings validate numerical inputs in the core and translate exceptions.
Returned arrays own their Python-visible storage. The allocator releases
the Python interpreter lock while computing. Result metadata includes
engine version, estimation method, normalization, window, annualization,
dataset hash and diagnostics. These fields describe reproducibility
inputs, not immutable historical data.

\subsection{Provider adapter and
limits}\label{provider-adapter-and-limits}

The Alpaca adapter loads credentials server-side from environment
variables or Streamlit secrets. Only market-data endpoints are called.
Historical requests explicitly select IEX or SIP and adjustment mode
\texttt{all}; pagination follows page tokens, rejecting repeated tokens
and excessive pagination. Authentication and malformed responses fail
clearly. Transport errors, HTTP 429 and server failures receive at most
four attempts with bounded waits. There is no silent feed fallback.
\href{https://docs.alpaca.markets/us/docs/about-market-data-api}{Alpaca's
access documentation} distinguishes coverage and subscription
requirements.

The adapter has a bounded, 15-minute in-memory history cache, isolated
per dashboard session and credential identity. Cache keys include symbols,
range, feed and adjustment mode. Requests end at least 15 minutes before
analysis time rather than in the future. On expiry or manual refresh the
entire adjusted window is fetched, and snapshots are compared for added,
removed and revised observations, including older dates. Returned histories
are defensive copies. This audits incremental changes without mixing
adjustment vintages; it does not reduce network volume. An overlap-only
refresh cannot guarantee detection of arbitrary revisions outside its overlap.

\href{https://docs.alpaca.markets/us/reference/stockbars}{Historical
bars documentation} lists adjustment choices, but an option name alone
does not establish dividend reinvestment economics. Known split/dividend
cases, account entitlement and display/redistribution permissions remain
unverified. Checked-in fixtures are synthetic. Mocked HTTP tests
establish software behavior, not successful live access. For an independently
specified corporate event, the audit helper compares the adjusted-price return
with the holder's wealth return
\[
r_{\mathrm{event}}=\frac{kP_{\mathrm{after}}+d}{P_{\mathrm{before}}}-1,
\]
where \(k\) is new shares per old share and \(d\) is cash per old share.
\begin{itemize}
\item \(r_{\mathrm{event}}\): Holder wealth return across the event; event is a label.
\item \(k,d\): New shares per old share, and cash distribution per old share.
\item \(P_{\mathrm{before}},P_{\mathrm{after}}\): Share prices before and after the event.
\end{itemize}

A two-for-one split from 100 to 50 gives zero wealth return. A cash dividend
of one with prices 100 and 99 also gives zero. Backward price adjustment
is not generally identical to cash-wealth arithmetic when the ex-date price
also changes for other reasons. Tests intentionally detect that discrepancy;
the helper cannot certify a provider convention without checked live events.

\subsection{Calendar and observation
policy}\label{calendar-and-observation-policy}

Only XNYS sessions whose close plus a 15-minute publication allowance
precedes the analysis time enter daily-return preparation. Explicit UTC
timestamps are converted to New York session dates. Prices must be
positive and finite; duplicate observations fail. Missing history at the
latest completed session fails. The selected universe is USD US stocks
and unleveraged ETFs; user confirmation is required, while automated
asset eligibility validation remains pending.

Prices are reindexed to the expected calendar before returns are
computed. Each return requires prices on adjacent expected sessions, so
a missing Monday cannot turn Friday-to-Tuesday movement into a daily
return. Only then are common complete return rows retained. Effective
row counts, excluded dates and a SHA-256 identity of the prepared
returns are recorded. No forward-filling occurs. Tests cover gaps,
weekends, incomplete sessions, pagination, retries, staleness and
estimator normalization.

Current shares use separately retrieved, unadjusted latest trades. A
price older than 15 minutes during an open session, or older than 15
minutes before the last completed close while closed, is flagged stale.
Prices more than 30 seconds into the future and nonpositive prices are rejected. Stale share valuation stops
the interface. A last trade is an observation, not an executable quote;
this threshold is an explicit policy, not a provider guarantee. Dollar
and hypothetical-weight inputs do not need the valuation request.

\subsection{Interface and deployment
boundary}\label{interface-and-deployment-boundary}

The Streamlit app supports a credential-free synthetic example,
attribution, exact versus local purchase effects, contribution
allocation and diagnostics. Market mode adds dollar/share/weight input,
feed selection and comparisons over sample and Ledoit--Wolf estimates at
63, 126 and 252 observations when available. Zero-volatility gradients
and infeasible caps are explained explicitly.

On 9 October 2026, the private Streamlit deployment successfully imported
and ran the C++ engine. Its synthetic workflow reproduced 13.74\% volatility
and the 200-dollar purchase of B. Authenticated IEX retrieval returned 379
common completed daily observations for SPY/BND without dropped return dates.
The hypothetical holdings of 600/400 dollars gave 8.4781\% estimated annual
volatility using 252 observations; allocation of a 200-dollar contribution
gave 7.4276\%. All six window/estimator comparisons converged. The Share
dialog showed the public-access option unchecked. These are hosted integration
checks at source revision \texttt{15ccd273}, separate from Ubuntu tests.
Credentials stayed in server secrets; no brokerage execution was performed.

\section{5. Numerical verification and
reproducibility}\label{numerical-verification-and-reproducibility}

\subsection{Evidence and reproduction}\label{evidence-and-reproduction}

\textbf{Symbols used in this subsection}
\begin{itemize}
\item \(w,w^*,v\): Feasible allocation, optimal allocation, and linear-oracle minimizing allocation; $v$ here is a vector, not variance.
\item \(f,\nabla f,{}^T\): Convex allocation objective, its gradient, and transpose.
\end{itemize}


At revision \texttt{fb1d449b},
\href{https://github.com/eshangark/portfolio-risk-engine/actions/runs/38088543928}{run 38088543928}
passed all 44 tests on Python 3.10 and 3.12, native sanitizers and all report
commands. The added audit checks 48 randomized risk instances, 48 allocation
solves against OSQP and an independent linear-oracle convex certificate,
covariance/dollar scales from $10^{-200}$ to $10^{200}$, extreme score differences,
negative risk contributions and nullspace portfolios. It also checks exact
historical-prefix invariance and independently known cash-flow ledgers.
For a feasible $w$ and a linear-oracle minimizer $v$ of $\nabla f(w)^Tv$,
convexity yields
\[0\le f(w)-f(w^*)\le \nabla f(w)^T(w-v).\]
The oracle fills the available capacity in increasing gradient order; it
does not reuse the projection routine. These bounded tests support the
implemented behavior, not a proof of reliability for every floating-point input.

Install a C++17 compiler, CMake and Eigen3, then run from the repository
root:

\begin{Shaded}
\begin{Highlighting}[]
\ExtensionTok{python} \AttributeTok{{-}m}\NormalTok{ pip install }\StringTok{\textquotesingle{}.[test,app]\textquotesingle{}}
\ExtensionTok{python} \AttributeTok{{-}m}\NormalTok{ unittest discover }\AttributeTok{{-}s}\NormalTok{ tests }\AttributeTok{{-}v}
\ExtensionTok{python} \AttributeTok{{-}m}\NormalTok{ examples.two\_asset }\AttributeTok{{-}{-}output}\NormalTok{ work/two{-}asset.json}
\ExtensionTok{python} \AttributeTok{{-}m}\NormalTok{ examples.engine\_report }\AttributeTok{{-}{-}output}\NormalTok{ work/engine{-}report.json}
\end{Highlighting}
\end{Shaded}

The
\href{https://github.com/eshangark/portfolio-risk-engine/blob/main/examples/engine_report.py}{C++
report script} records synthetic data identity, source hashes, runtime
environment, derivative errors and allocation diagnostics. The Python
3.12.15 / NumPy 2.5.3 run measured maximum weight-gradient error
\(2.78\cdot10^{-17}\) and dollar-gradient error \(4.07\cdot10^{-20}\)
against the independent reference. The worked allocation had zero
reported budget, bound and projected-gradient residuals. These values
describe this fixture, not worst-case error bounds. Each Python job also
stores JSON artifacts. Source hashes identify the evaluated code; later
documentation edits do not change the historical record.

{\def\LTcaptype{none} % do not increment counter
\begin{longtable}[]{@{}
  >{\raggedright\arraybackslash}p{(\linewidth - 2\tabcolsep) * \real{0.5000}}
  >{\raggedright\arraybackslash}p{(\linewidth - 2\tabcolsep) * \real{0.5000}}@{}}
\toprule\noalign{}
\begin{minipage}[b]{\linewidth}\raggedright
Test family
\end{minipage} & \begin{minipage}[b]{\linewidth}\raggedright
Independent evidence
\end{minipage} \\
\midrule\noalign{}
\endhead
\bottomrule\noalign{}
\endlastfoot
Risk and attribution & Hand variance .01888; NumPy formula; Euler
identity \\
Gradients & Analytical, AD and central-difference comparisons \\
Dollar normalization & Recomputed holdings and feasible transfer
derivatives \\
Tape semantics & Shared nodes, repeated passes, lifetime, reset and
domains \\
Projection & Budget and bounds, idempotence, variational inequality \\
Allocation & OSQP objectives for full-rank and singular matrices \\
Degenerate cases & Zero budget/covariance, iteration limit, infeasible
caps \\
Covariance & PSD validation, rank deficiency, scaling and permutation \\
Estimation & Hand covariance and reconstructed fitted shrinkage \\
Data & Mock pagination/retries, calendar gaps, staleness and caching \\
Interface & Synthetic display, infeasible caps, universe confirmation \\
\end{longtable}
}

Native tests additionally run with AddressSanitizer and
UndefinedBehaviorSanitizer. Passing this finite set provides evidence
for covered cases, not a proof of memory safety for every input. Tests
use synthetic inputs without provider credentials.

\subsection{Tolerances and
interpretation}\label{tolerances-and-interpretation}

\textbf{Symbols used in this subsection}
\begin{itemize}
\item \(n,\epsilon_{mach},S_{ij}\): Asset count, machine epsilon, and covariance entry at row $i$, column $j$.
\item \(h,O(h^2),O(\epsilon_{mach}/h)\): Finite-difference step (not holdings here), truncation-error order, and cancellation-error order.
\item \(\max,|\cdot|\): Maximum operator and absolute value.
\end{itemize}


Calculations use float64. Covariance validation tolerates asymmetry and
negative eigenvalues only at \(100n\epsilon_{mach}\max|S_{ij}|\). It
symmetrizes and clips negative eigenvalues only inside that allowance,
then uses the same effective matrix for risk and gradients. Materially
indefinite matrices are rejected. This relative rule has no absolute
floor hiding errors at small scales. The risk API reports the Frobenius
norm of its covariance correction.

Weights must sum to one within \(10^{-12}\); they are not silently
renormalized. Dollar normalization is explicit. Central finite
differences have truncation error \(O(h^2)\) and cancellation error
roughly \(O(\epsilon_{mach}/h)\). Multiple step sizes are used, with
relative tolerances up to \(3\cdot10^{-6}\) for coarse steps and
absolute allowances near zero. See
\href{https://docs.pytorch.org/docs/2.14/notes/gradcheck.html}{gradcheck
mechanics}; PyTorch is not a dependency.

OSQP checks request absolute and relative tolerances \(10^{-10}\).
Comparisons allow objective relative error \(10^{-7}\) and absolute
error \(10^{-10}\), alongside engine convergence and feasibility
diagnostics. Singular problems may have several minimizers, so equal
weights are not required. Near-zero volatility can make sensitivities
ill-conditioned even when variance evaluation is stable. Extreme
floating-point scale testing remains a hardening task.

\subsection{Worked interpretation}\label{worked-interpretation}

\textbf{Symbols used in this subsection}
\begin{itemize}
\item \(g_j,V,j\): Volatility-gradient entry, total portfolio value, and asset index.
\end{itemize}


The synthetic portfolio has volatility .1374045123 and contributions
.1152800569 and .0221244554. Fixed-alpha shrinkage changes volatility
from .1374045 at zero to .1140175 at one; this spread measures model
dependence, not statistical confidence. The purchase example shows
increasing error in a local derivative approximation as purchases grow.

For the same covariance, holdings \((600,400)\), contribution 200 and
caps \((.6,.6)\), the solution purchases \((0,200)\) and ends at
\((.5,.5)\) with variance .0155. Along the feasible interval the
variance derivative favors additional B, but the no-sale bound prevents
reducing A below .5. The optimum is at a boundary despite a nonzero
unconstrained gradient.

\subsection{Historical evaluation}\label{remaining-empirical-work}
\subsection*{Historical contribution backtest}
\begin{itemize}
\item \(h_{t-1},h_t^-,h_t,x_t\): Prior end-of-day holdings, holdings after deposits before returns, current end-of-day holdings, and current purchase dollars; minus labels a pre-return state.
\item \(t,r_t,R_t\): Day index, vector of asset returns, and portfolio return excluding the direct deposit effect.
\item \(\mathbf1,{}^T,\odot\): All-ones vector, transpose, and entrywise multiplication.
\item \(I_t,I_T,I_0,\prod_t\): Time-weighted wealth index, its final value, initial index of one, and product over days.
\item \(T,s,I_s,\min_t,\max\): Final observation index, earlier day index in the running maximum, index at that earlier day, minimum over days, and maximum operator.
\item \(s(R_t)\): Sample standard deviation of the daily portfolio-return series; $s$ is a function name here, not the earlier covariance scale.
\end{itemize}

The cumulative evaluator is separate from the fixed-holdings variance study.
At each decision date it fits Ledoit--Wolf covariance on the preceding 126
observations, before consuming that day's return. Initial holdings are 1,000
dollars divided equally; every 21 observations adds 200 dollars. There are
no sales and no concentration cap below 100\%. The comparator makes equal-dollar
purchases on the same dates with the same cash flows. Neither policy chooses
parameters from future outcomes.

Let $h_t^-=h_{t-1}+x_t$ be holdings after any contribution but before the
daily return. Then $h_t=h_t^-\odot(\mathbf1+r_t)$ and the portfolio return is
$R_t=\mathbf1^Th_t/\mathbf1^Th_t^- -1$. The time-weighted index
$I_T=\prod_t(1+R_t)$ removes the direct deposit effect. Drawdown is
$\min_t\{I_t/\max_{0\le s\le t}I_s-1\}$ with $I_0=1$; annualized
volatility is $\sqrt{252}\,s(R_t)$. Terminal wealth is reported separately
because it includes deposits. Holdings drift daily; this differs from the
constant-weight portfolio used in variance calibration.

Zero returns must give terminal wealth equal to initial wealth plus deposits,
zero time-weighted return and zero drawdown. A common 1\% daily return has
a closed-form wealth ledger, independently tested. Changing observations
after a cutoff must leave the earlier cash-flow ledger and decisions exactly
unchanged. Contiguous input arrays fix layout-dependent reduction differences
without weakening that exact-prefix test. Missing market return dates suppress
the hosted backtest, since skipping them would omit unknown wealth changes.
Fees, taxes, spreads, executable prices, point-in-time vintages and survivorship
effects are not modeled.

The seeded four-regime Gaussian stress simulation contains 1,008 observations
and four assets. It evaluates 882 observations and 42 contributions totaling
8,400 dollars. On the Python 3.10 runner, optimized / equal-purchase terminal
wealth was 8,460.41 / 8,507.02 dollars, time-weighted return 1.3306\% / 2.3450\%,
volatility 23.3873\% / 23.2907\%, and drawdown $-38.0320\%$ / $-37.0318\%$.
The optimized policy was worse on these realized measures: minimizing fitted
variance at each contribution is not a guarantee of lower subsequent variance,
drawdown or higher return. This is synthetic stress evidence, not market performance.

The extended hosted market-data run used SPY, BND, GLD and QQQ with 1,513
complete IEX-adjusted daily returns, retrieved on 9 October 2026 at
22:05:14 UTC under revision \texttt{cb737b9b}. The backtest evaluated 1,387
observations from 5 April 2021 through 9 October 2026, making 67 contributions
totaling 13,400 dollars. All solves converged; no missing daily returns were
skipped. Aggregate results were:
\begin{center}
\begin{tabular}{lrr}
\toprule
Measure & Minimum variance & Equal purchases \\
\midrule
Terminal wealth (dollars) & 17,738.59 & 23,078.07 \\
Time-weighted return & 31.99\% & 93.75\% \\
Annualized volatility & 7.20\% & 12.27\% \\
Maximum drawdown & $-18.66\%$ & $-21.70\%$ \\
\bottomrule
\end{tabular}
\end{center}
Lower volatility and drawdown accompanied lower return and wealth on this
chosen history. This tradeoff does not establish general superiority.
The independent fixed-holdings study on the same snapshot completed 360
model evaluations across 60 test blocks. Mean calibration ratios for sample
/ Ledoit--Wolf were 1.1684 / 1.1366 (63 observations), 1.1975 / 1.1765 (126),
and 1.1870 / 1.1696 (252). Ratios above one indicate underestimation of
subsequent variance on average in these groups. Present-day adjustments,
surviving selected tickers and the unverified provider dividend convention
prevent treating this as an unbiased investment backtest. These aggregates describe a recorded historical snapshot, not a fresh-data guarantee.
The \href{https://github.com/eshangark/portfolio-risk-engine/blob/main/examples/live_report.py}{live report script}
reproduces the evaluation protocol against retrieved history; raw market histories
and credentials are not committed.

\subsection*{Variance calibration and limitations}
\begin{itemize}
\item \(w_t,r_s,t,s\): Weights fitted at decision time $t$, asset returns in later test observation $s$, decision index, and test-observation index.
\item \(\operatorname{Var}_{\mathrm{test}},{}^T\): Sample variance within the held-out test block and transpose.
\item \(L^1\): Sum of absolute component changes in consecutive allocation vectors; the superscript identifies the norm, not a power.
\end{itemize}


Unperformed work includes live corporate-action validation, SIP entitlement checks
and broader empirical analysis. The implemented
research protocol fits windows of 63, 126 and 252 observations through a cutoff,
then evaluates the next 21 observations, using nonoverlapping test blocks.
Each cutoff poses an independent contribution problem with fixed initial
holdings of 100 dollars per asset and a 200-dollar contribution. It is not
a cumulative trading simulation. For fitted weights \(w_t\), it compares
predicted annual variance with \(252\operatorname{Var}_{\mathrm{test}}(r_s^Tw_t)\),
using sample normalization within the test block. The calibration ratio and
successive allocation \(L^1\) changes summarize variance fit and stability.
Future-perturbation tests ensure that changing test observations cannot change
the first fitted allocation. Present-day adjusted history and a surviving-symbol
universe can still introduce vintage or selection bias despite correct slicing.

The hosted runner at revision \texttt{8da7a371} completed 36 evaluations across
six test blocks on the SPY/BND snapshot retrieved at
\texttt{2026-10-09T21:38:30.377874+00:00}. Mean calibration ratios were:
\begin{center}
\begin{tabular}{rrr}
\toprule
Training observations & Sample & Ledoit--Wolf \\
\midrule
63 & 0.9040 & 0.8344 \\
126 & 1.0251 & 0.9666 \\
252 & 1.2418 & 1.1551 \\
\bottomrule
\end{tabular}
\end{center}
Each group contains only six evaluations. Mean successive allocation
$L^1$ change was zero in each group, consistent with unchanged boundary
solutions. This selected two-asset example cannot establish estimator superiority
or general allocation stability. The daily-return dataset SHA-256 identity is
\begin{quote}\ttfamily\small
2a3c00006e0e1b8b35309c375b0c00989\\
d291c180d78ca569382c371e9905080
\end{quote}

The benchmark compares validated risk APIs against the vectorized NumPy
reference, with five warm-up calls and seven batches of 30 calls at 2, 10
and 50 assets. It records median and percentile timings, compiler, platform,
NumPy build configuration and thread settings. These are end-to-end costs
including validation, eigendecomposition and Python/C++ crossing, not isolated
kernel timings. Isolating pure binding overhead remains a separate benchmark.
API timings do not imply any investment-return advantage.

\subsection{Measured timings}

The synthetic factor study uses seed 24109, 504 observations and three
assets. It completed 72 evaluations with mean calibration ratio 1.0096.
This is a reproducible exercise of the protocol under synthetic data,
not empirical evidence about market covariance or an estimated confidence
level. Full evaluation records are retained in the workflow artifacts.

On the Python 3.12 Ubuntu x86-64 runner, using GCC 13.3 and one configured
BLAS/OpenMP thread, measured median microseconds per validated API call were:

\begin{center}
\begin{tabular}{rrrr}
\toprule
Assets & NumPy & C++ analytical & C++ AD \\
\midrule
2 & 34.44 & 3.48 & 4.48 \\
10 & 45.14 & 9.98 & 30.41 \\
50 & 285.10 & 182.34 & 884.69 \\
\bottomrule
\end{tabular}
\end{center}

The analytical C++ path was faster for these cases, while the AD path was
slower than NumPy at 50 assets. Graph construction and traversal have a
material cost. This result depends on the tested APIs, runner and dimensions;
it is not a claim that C++ or AD is universally faster. The
\href{https://github.com/eshangark/portfolio-risk-engine/blob/main/examples/benchmark.py}{benchmark script}
records percentile timings and the numerical-library environment.


\subsection{Comprehension questions}\label{comprehension-questions}

\textbf{Symbols used in this subsection}
\begin{itemize}
\item \(g_j,V,j\): Volatility-gradient entry, total portfolio value, and asset index.
\end{itemize}


\begin{enumerate}
\def\labelenumi{\arabic{enumi}.}
\tightlist
\item
  Why does a 60\% capital holding contribute about 84\% of volatility?
\item
  Why is a new-dollar derivative \((g_j-\sigma)/V\) instead of
  \(g_j/V\)?
\item
  Why can singular covariance be valid while optimal weights are
  nonunique?
\item
  Why can a constrained optimum have a nonzero ordinary gradient?
\item
  Which assumptions justify square-root-of-time annualization?
\end{enumerate}

\section{6. References and glossary}\label{references-and-glossary}

Sources consulted on 2026-10-08. Documentation is mutable; recheck when
implementing integrations. Entries distinguish material actually
consulted from full proofs not yet reviewed.

\begin{enumerate}
\def\labelenumi{\arabic{enumi}.}
\tightlist
\item
  \textbf{MOSEK, Portfolio Optimization Cookbook},
  \href{https://docs.mosek.com/portfolio-cookbook/markowitz.html}{chapter
  2}, and
  \href{https://docs.mosek.com/portfolio-cookbook/risk_parity.html}{chapter
  9}. Consulted portfolio formulation and Euler attribution. Our
  contribution-only allocation is a project-specific specialization.
\item
  \textbf{Olivier Ledoit and Michael Wolf (2004)}, ``A well-conditioned
  estimator for large-dimensional covariance matrices,'' \emph{Journal
  of Multivariate Analysis} 88(2), 365--411,
  \href{https://www.sciencedirect.com/science/article/pii/S0047259X03000964}{DOI
  10.1016/S0047-259X(03)00096-4}. Publisher abstract and indexed
  author-paper introduction consulted for motivation and scope of
  asymptotic covariance-loss claims. Full proof review remains pending;
  author-hosted PDF retrieval was unsuccessful. No theorem from the
  unread proof is claimed as newly verified.
\item
  \textbf{scikit-learn},
  \href{https://scikit-learn.org/stable/modules/generated/sklearn.covariance.LedoitWolf.html}{LedoitWolf
  API}. Consulted the scaled-identity formula and centering contract.
  The estimator wrapper is implemented and empirical normalization is
  tested.
\item
  \textbf{Atılım Güneş Baydin, Barak A. Pearlmutter, Alexey Andreyevich
  Radul, Jeffrey Mark Siskind (2018)}, ``Automatic Differentiation in
  Machine Learning: a Survey,'' \emph{JMLR} 18(153), 1--43.
  \href{https://www.jmlr.org/beta/papers/v18/17-468.html}{Publisher
  record},
  \href{https://jmlr.org/papers/volume18/17-468/17-468.pdf}{paper}.
  Consulted section 3, reverse traversal and accumulation. Its title
  does not imply machine learning is part of our project.
\item
  \textbf{Neal Parikh and Stephen Boyd (2014)}, ``Proximal Algorithms,''
  \emph{Foundations and Trends in Optimization} 1(3), 127--239.
  \href{https://web.stanford.edu/~boyd/papers/pdf/prox_algs.pdf}{Author-hosted
  paper}. Consulted projected gradient as proximal gradient for a set
  indicator; project-specific capped-simplex derivation is provided in
  chapter 3.
\item
  \textbf{PyTorch},
  \href{https://docs.pytorch.org/docs/2.14/notes/gradcheck.html}{Gradcheck
  mechanics}. Consulted central finite differences and analytical
  comparison. Reference tests do not use PyTorch.
\item
  \textbf{pybind11},
  \href{https://pybind11.readthedocs.io/en/stable/}{official
  documentation}. Consulted C++/Python integration overview. The
  extension uses CMake and scikit-build-core; tests exercise arrays and
  handle lifetime.
\item
  \textbf{Alpaca},
  \href{https://docs.alpaca.markets/us/docs/about-market-data-api}{Market
  Data API overview} and
  \href{https://docs.alpaca.markets/us/reference/stockbars}{historical
  bars endpoint}. Consulted feed/access distinctions, pagination and
  adjustment options. Data entitlement and display/redistribution terms
  are not yet validated.
\item
  \textbf{Streamlit},
  \href{https://docs.streamlit.io/deploy/streamlit-community-cloud}{Community
  Cloud},
  \href{https://docs.streamlit.io/deploy/streamlit-community-cloud/share-your-app}{privacy/sharing},
  and
  \href{https://docs.streamlit.io/deploy/streamlit-community-cloud/deploy-your-app/app-dependencies}{dependencies}.
  Consulted private-app allowance and Linux dependency mechanism. Private
  hosted build and IEX integration were checked on 9 October 2026.
\end{enumerate}

\clearpage
\subsection{Glossary}\label{glossary}

{\def\LTcaptype{none} % do not increment counter
\begin{longtable}[]{@{}
  >{\raggedright\arraybackslash}p{(\linewidth - 2\tabcolsep) * \real{0.5000}}
  >{\raggedright\arraybackslash}p{(\linewidth - 2\tabcolsep) * \real{0.5000}}@{}}
\toprule\noalign{}
\begin{minipage}[b]{\linewidth}\raggedright
Term
\end{minipage} & \begin{minipage}[b]{\linewidth}\raggedright
Meaning here
\end{minipage} \\
\midrule\noalign{}
\endhead
\bottomrule\noalign{}
\endlastfoot
Covariance & Expected product of centered returns \\
PSD & Every quadratic form is nonnegative \\
Eigenvalue & Variance scaling along an eigenvector direction \\
Conditioning & Sensitivity of a numerical problem to input
perturbations \\
Shrinkage & Blend an estimate with a structured target \\
Gradient & Vector of unconstrained coordinate derivatives \\
Directional derivative & Local change along a specified direction \\
Adjoint & Derivative of the output with respect to an intermediate
node \\
Tape & Stored computational graph and intermediate values \\
Projection & Nearest feasible point in the chosen norm \\
KKT conditions & Feasibility, stationarity, multiplier signs and
complementarity \\
Residual & Quantified violation of a constraint or optimality
condition \\
Annualization & Conversion under declared time-aggregation
assumptions \\
Sensitivity & Change under specified perturbations, not a confidence
interval \\
\end{longtable}
}

\end{document}
